% Weekly exercise sheet for Lecture 9; also included in the synthesis sheet.
\ifdefined\ExercisesSynthesis
\else
\documentclass[11pt]{article}
\usepackage[utf8]{inputenc}
\usepackage[T1]{fontenc}
\usepackage[english]{babel}
\usepackage{lmodern,microtype}
\usepackage[margin=0.9in]{geometry}
\usepackage{amsmath,amssymb,amsthm,mathtools}
\usepackage{enumitem,hyperref,needspace}
\newcommand{\E}{\mathbb E}
\newcommand{\Qbb}{\mathbb Q}
\newcommand{\dd}{\,\mathrm d}
\newcommand{\Pb}{\mathbb P}\newcommand{\F}{\mathcal F}\newcommand{\Var}{\operatorname{Var}}
\theoremstyle{definition}
\newtheorem{exercise}{Exercise}[section]
\newlist{exparts}{enumerate}{1}
\setlist[exparts,1]{label=\textbf{(\alph*)},itemsep=0.3em,topsep=0.4em,leftmargin=2em}
\newcommand{\lecturesummary}[1]{\begin{quote}\small\textbf{Focus.} #1\end{quote}}
\title{Stochastic Calculus and Option Pricing\\Lecture 9 Exercises}
\author{}
\date{}
\AddToHook{env/exercise/before}{\Needspace{7\baselineskip}}
\expandafter\newif\csname ifIncludeCore\endcsname\IncludeCoretrue
\expandafter\newif\csname ifIncludeFurther\endcsname\IncludeFurthertrue
\begin{document}
\maketitle
\vspace{-1.2em}
\setcounter{section}{8}
\fi

\section{Week 9 --- Vanilla Options in Practice}
\lecturesummary{Check contract units, sensitivity signs, gamma-weighted hedging outcomes, and executable surface-arbitrage tests.}

\ifIncludeCore

\lecturesummary{Core workload: about 90 minutes, including 10 minutes to check signs, assumptions, and conclusions. Suggested solving times: 9.1: 20 min; 9.2: 20 min; 9.3: 20 min; 9.4: 20 min. These are planning estimates after studying the lecture.}
\setcounter{exercise}{0}
\begin{exercise}[Variation --- A nonnegative call spread can have negative vega]
In Black--Scholes with no dividends, let
\(V=C(K_1)-C(K_2)\), where \(K_1<K_2\) and both calls have maturity
\(T\). At time \(t\), write \(\tau=T-t\). You may use
\(\mathcal V(K)=S\phi(d_1(K))\sqrt\tau\).
\begin{exparts}
\item Show that the spread vega is positive exactly when \(|d_1(K_1)|<|d_1(K_2)|\). Explain why positivity of each individual call vega does not settle the sign of their difference.
\item Find the spot level at which spread vega is zero. Determine its sign on either side of that level.
\item Use \(\mathcal V=S^2\sigma\tau\Gamma\) to relate the signs of spread gamma and spread vega. What does this say about curvature of a bounded spread payoff?
\item Explain why nonnegative terminal payoff and positive price do not imply that the price increases with volatility.
\end{exparts}
\end{exercise}

\begin{exercise}[Variation --- Equal realized variance, different hedge outcomes]
For a continuously delta-hedged short option, suppose the discounted
terminal hedging profit can be written
\[
\widetilde\Pi_T=-\frac12\int_0^T w_t
\bigl(\sigma_{\rm actual}(t)^2-\widehat\sigma^2\bigr)\dd t,
\qquad w_t=e^{-rt}S_t^2\Gamma_t\geq0.
\]
Consider two half-periods of equal length. In a controlled comparison,
take their weights to be constants \(w_1>w_2>0\).
\begin{exparts}
\item Compare variance errors \((+\delta,-\delta)\) and \((-\delta,+\delta)\), with \(\delta>0\) small enough to keep variances positive. Both errors integrate to zero. Compute the two profits.
\item Which part of the path matters most for the hedging result? Why is comparison of unweighted annualized realized and quoted volatility insufficient?
\item If the actual variance exceeds the hedge variance at every date, what sign follows for the short-option profit under the stated gamma assumption?
\item If the payoff is a call spread with changing gamma sign, which conclusion above can fail? Explain without changing the self-financing accounting.
\end{exparts}
\end{exercise}

\begin{exercise}[Variation --- Follow a maturity slice at fixed forward moneyness]
Assume constant \(r\), no dividends, and a true risk-neutral martingale
\(X_t=e^{-rt}S_t/S_0\). Fix log-forward moneyness
\(x=\log(K/(S_0e^{rT}))\), and define
\(c(T,x)=C(T,S_0e^{rT+x})/S_0\). Assume the price surface is smooth.
\begin{exparts}
\item Use the chain rule to find \(\partial_T c(T,x)\) at fixed \(x\). Why is this different from differentiating \(C(T,K)\) at fixed strike?
\item Write \(c(T,x)=\E^{\Qbb}[(X_T-e^x)^+]\). Use conditional Jensen at two maturities to show that this quantity is nondecreasing in \(T\). Deduce the derivative's sign without using a later local-volatility formula.
\item Under Black--Scholes, the normalized call at fixed \(x\) depends on total variance \(w=\sigma_{\rm imp}^2T\) and increases strictly with \(w>0\). Explain why decreasing annualized implied volatility can be consistent with the calendar restriction.
\item Compare 30\% implied volatility at \(T=1/4\) with 20\% at \(T=1\), at the same \(x\). What does this two-point check establish, and what does it leave untested?
\end{exparts}
\end{exercise}

\begin{exercise}[Practice --- Repair a quote when strikes are not equally spaced]
At one maturity, rates are zero, \(S_0=100\), and European call
quotes are \(C(90)=14\), \(C(100)=x\), and \(C(120)=3\). Ignore
bid--ask spreads for this exercise.
\begin{exparts}
\item Using monotonicity and call-spread slope bounds, find all restrictions on \(x\) from the two neighboring intervals.
\item Derive the unequal-strike convexity restriction. Also include the zero-strike call \(C(0)=S_0=100\) and convexity across strikes \(0,90,100\). Combine the restrictions to obtain the admissible interval for these quote checks.
\item If \(x=11\), construct a nonnegative-payoff weighted butterfly with negative initial cost. Explain how to turn it into a zero-cost arbitrage.
\item Why is the equally weighted package with holdings \((1,-2,1)\) an invalid convexity test for these strikes? Examine its payoff above 120.
\end{exparts}
\end{exercise}
\fi

\ifIncludeFurther
\Needspace{10\baselineskip}
\subsection*{Optional further practice}
\noindent These problems are supplementary; the core workload is complete after Exercises 9.1--9.4.
\setcounter{exercise}{4}
\begin{exercise}[A stock split changes units, not economics]
Suppose one old share is replaced by \(a>0\) new shares, so that the quoted
stock price and strike become \(S'=S/a\) and \(K'=K/a\).
\begin{exparts}
  \item Compare the old call payoff \((S_T-K)^+\) with the payoff on \(a\)
  new calls struck at \(K/a\). Explain why the option position must also be
  rescaled by \(a\).
  \item If \(C(S,K)\) is the old call value and \(C'(S',K')\) the value per
  new call, derive the relation between them by no-arbitrage.
  \item Use the chain rule to relate the old and new Delta and Gamma. Which
  sensitivity is invariant to the change of share denomination?
\end{exparts}
\end{exercise}
\fi

\ifdefined\ExercisesSynthesis\else
\end{document}
\fi
